\section{Reproduction}\label{sec:reproduction}
The derivations, checks and tests use only the Python standard library and
run offline on the frozen inputs; from the repository root:
\begin{footnotesize}
\begin{verbatim}
python3 -m unittest discover -s tests
python3 scripts/check_certification.py
python3 scripts/run_contracts.py
python3 -B experiments/reactivation_probes/derive.py
python3 -B experiments/revision_bundles/derive.py
python3 -B experiments/revision_bundles/addendum.py
python3 -B experiments/revision_bundles/addendum2.py
python3 scripts/paper_numbers.py
python3 scripts/supplement_tables.py
\end{verbatim}
\end{footnotesize}
\code{paper_numbers.py} writes \code{paper/numbers.tex}, from which every number
in the paper is typeset, and asserts identities among the derived files;
\code{supplement_tables.py} writes this document's
tables and number macros from the same files. Re-executing the probes or the
collection needs the pinned servers or network access and produces new IDs
and times; the derivations then compare classifications, not bytes.

Each directory under \code{figure-artifacts/} and \code{table-artifacts/}
holds its frozen inputs, derivation, drawing or table code, exact plotted
values, caption, metadata, outputs and a SHA-256 manifest, and rebuilds alone.
Figures 1 and 2 are TikZ and PGFPlots code generated by \code{plot.py} from
\code{plot-data.csv} and typeset with XeLaTeX in the acmart typefaces at the
column width; \code{plot.py} checks page width, font embedding, type sizes and
clipping, and rebuilds are byte-identical. This supplement is built with
XeLaTeX from TeX Live 2023:
\begin{footnotesize}
\begin{verbatim}
python3 scripts/build_figure_artifacts.py
python3 scripts/build_supplement.py
\end{verbatim}
\end{footnotesize}

The artifact that accompanies the submission holds the code, the frozen data,
the figure and table directories, this supplement with its sources, and
\code{paper/numbers.tex} with every number printed in the paper. It was
checked in a fresh copy before release: the tests pass,
\code{paper_numbers.py} and \code{supplement_tables.py} reproduce their
outputs unchanged, and \code{build_supplement.py} rebuilds this document byte
for byte.

Each kind of evidence certifies a bounded statement. The contract checker
certifies properties of the supplied graph and scope; the certification code
decides the stated models; local executions show what one pinned
implementation did in one run; documentation states intended service
behavior; and the replay evaluates certificates against reconstructed edit
histories of one wiki. None of them measures how often a deployed archive
makes the inference the paper examines.
